\documentclass[10pt,a4paper]{amsart}
\usepackage[margin=19mm]{geometry}
\usepackage{amsmath,amssymb,mathtools}
\usepackage[scr=rsfs,cal=euler]{mathalfa}
\usepackage{xfrac}
\usepackage[unicode,hidelinks,bookmarks=false]{hyperref}
\newcommand{\ie}{i.e.\ }
\newcommand{\cf}{cf.\ }
\newcommand{\resp}{respectively\ }
\newcommand{\Subsection}{\subsection}
\providecommand{\nobreakdash}{\nobreak\hbox{-}}
\setlength{\emergencystretch}{2em}
\multlinegap0pt
\usepackage{amssymb}
\usepackage{enumerate}
\usepackage{xcolor}
\newtheorem{theorem}{Theorem}
\newcommand{\Z}{\mathbb{Z}}
\newcommand{\si}{\mathcal{S}}
\title{Essential quotients of the cohomology of products of symmetric and alternating groups}
\author{Dana Hunter, Dev Sinha}
\date{Source: arXiv:2609.14781v1 (CC BY 4.0)}
\begin{document}
\maketitle

\noindent\emph{Source and licence.} Essential quotients of the cohomology of products of symmetric and alternating groups. Version arXiv:2609.14781v1 (CC BY 4.0), distributed by arXiv under the Creative Commons Attribution 4.0 International licence, http://creativecommons.org/licenses/by/4.0/, as recorded in the arXiv licence metadata for this exact version and shown on the abstract page https://arxiv.org/abs/2609.14781v1. Full licence notice: \texttt{licenses/arxiv-2609.14781-CC-BY-4.0.txt}.\par\smallskip

\noindent\textit{Editorial scope.} Counted unit: the article's theorem Wninvts (source0.tex lines 615--620). The article's complete original proof of it is delivered with it (source0.tex lines 622--673, one \texttt{proof} environment of 52 lines). No mathematics is written, repaired or re-derived: the statement and the proof are the article's own lines, in the article's own notation, and no retained line is substituted or re-flowed.  No further source-local context is needed: every cross-reference of every retained line resolves inside the two delivered spans.  Nothing was omitted from any retained span, so the package carries no omission record.  No retained line contains a citation, so the wrapper carries no bibliography and no bibliography entry.  No retained line contains a figure environment, an image-inclusion command or drawing code, so no image is delivered and the package declares no asset dependency.  Editorial formatting only, in addition: an AMS \texttt{amsart} preamble that restates the article's own declarations and macros used by the retained lines, namely the environments \texttt{theorem} on separate counters, together with the standard packages those lines need.  The retained environments print in this document as Theorem 1 in order of appearance, whereas the article prints the same items with the numbering of its own full document.  The complete native source of the article as distributed in the arXiv e-print for this exact version is bundled unchanged as \texttt{sources/210356/source0.tex}.
\begin{theorem}\label{Wninvts}
    For $n>1$, the quotient of $H^*(E_{alt})^{W_n}$, 
    by its intersection with the negligible ideal
in $H^*(E_{alt})$  is spanned
    by the image of the subring $\left( \Z/2[b]^{\otimes n}\right)^{\si_n}$ of $ H^*(E_{alt})^{W_n}$.
\end{theorem}

\begin{proof}
 %Let $\sigma$ denote the automorphism of $\Z/2[x_1, x_2]$ which exchanges $x_1$ and $x_2$.
    %It thus fixes
    %$b$ while exchanging $c_+$ with $c_-$.
    %Let $\tau_{i,j} \in W_n$ denote the automorphism of 
    %$\Z/2[x_1, x_2]^{\otimes n}$ which applies $\sigma$ on
    %both $i$th and $j$th factors.  
    We first treat the $n=2$ case.  Let $H \subset W_2$ be the subgroup generated by $\tau_{1,2}$ and the permutation of the two factors of $V_2$.   
    The action of $W_2$, and thus that of $H$, permutes the basis of tensor products of monomials.  We may therefore decompose a $W_2$-invariant $x$ as a sum 
    $$x = \sum_i \sum_{(\sigma) \in H/K_i} \sigma (f_i \otimes g_i),$$
    where $f_i$ and $g_i$ are monomials and 
    each $K_i$ is a subgroup of $H$.
      Explicitly, considering possible subgroups $K_i \subset H \cong C_2 \times C_2$, each sum
      $\sum_{(\sigma) \in H/K_i} \sigma (f_i \otimes g_i)$ is one of the following, where $f$, $f_+$, $f_-$, $g$, $g_+$ and $g_-$ are all monomials:
\begin{enumerate}
    \item $f \otimes f$ fixed by $\tau$;
    \item $f_+ \otimes f_+ + f_- \otimes f_-$, where $\tau(f_+) = f_-$ and $f_+ \neq f_-$;
    \item $f_+ \otimes f_- + f_- \otimes f_+$, where $\tau(f_+) = f_-$ and $f_+ \neq f_-$;
    \item $f \otimes g + g \otimes f$, where $f$ and $g$ are fixed by $\tau$;
    \item $f_+ \otimes g_+ + g_+ \otimes f_+ + f_- \otimes g_- 
    + g_- \otimes f_-$, where  $\tau(f_+) = f_-$, $f_+ \neq f_-$,$\tau(g_+) = g_-$ and $g_+ \neq g_-$.
\end{enumerate}
The invariants of $\Z/2[b,c_+,c_-]$ under $\tau$ are $\Z/2[b,c]$, where $c = c_+ + c_-$.  But  $c_+ +  c_- = {x_1}^2 x_2 + x_1 {x_2}^2$, which is negligible.  Thus any $\tau$ invariant is equivalent to an element of $\Z/2[b]$.  
Also, because $c_+ \sim c_-$  any monomial $f_+ \in \Z/2[b, c_+, c_-]$ will be equivalent to $\tau(f_+)$ modulo the negligible ideal.  Considering the orbit types above in the quotient by the negligible ideal we deduce
the following equivalences:
\begin{enumerate}
    \item $f \otimes f$ is equivalent to an element of $(\Z/2[b]^{\otimes 2})^{\si_2}$;
    \item $f_+ \otimes f_+ + f_- \otimes f_-$ is zero  because $f_+ \sim f_-$;
    \item $f_+ \otimes f_- + f_- \otimes f_+$ is zero;
    \item $f \otimes g + g \otimes f$ is equivalent to an element of $(\Z/2[b]^{\otimes 2})^{\si_2}$;
    \item $f_+ \otimes g_+ + g_+ \otimes f_+ + f_- \otimes g_- 
    + g_- \otimes f_-$ is zero.
\end{enumerate}
Thus, considering the $H$-orbit monomial decomposition of any $W_2$ invariant, it will be equivalent to an element of $(\Z/2[b]^{\otimes 2})^{\si_2}$, establishing the result in this case.


In general, let $H_{i,j}$ be the subgroup of $W_n$ generated by
$\tau_{ij}$ and the permutation of the $i$th and $j$th
factors of $V_2$.
%Let $\sum_k f_{k,1} \otimes \cdots \otimes f_{k,n} \in \Z/2[b, c_+, c_-]^{\otimes n}$ be a $W_n$ invariant, where each $f_{k,n}$ is a monomial. 
Decompose a $W_n$ invariant into $H_{i,j}$ orbits of tensor products of monomials,
each consisting of a sum of one, two or four tensor products of monomials as in the argument above.
Following the orbit analysis above, 
each such orbit will be equivalent to sum of one or two 
tensor products of monomials
whose $i$th and $j$th entries are in $\Z/2[b]$ and which
is invariant with respect to permuting the $i$th and $j$th entries.
    Applying this argument
    for all $i, j$ we see that modulo
    negligible elements, any $W_n$ invariant is equivalent to
    an element of $\left(\Z/2[b]^{\otimes n}\right)^{\si_n}$.
\end{proof} 

\end{document}
